\chapter*{Innovations}
\addcontentsline{toc}{chapter}{Innovations}

This chapter describes the practical applications and innovation potential of the research results presented in this thesis.

\section*{Industrial Applications of Research Results}

\begin{enumerate}
    \item \textbf{Real-time reliability monitoring for LLM deployments.} The uncertainty estimation pipeline developed in Study~1 enables automatic detection of unreliable LLM outputs at inference time, requiring only a single forward pass with no additional models or computational overhead. This is directly applicable to high-stakes domains:
    \begin{itemize}
        \item \textit{Healthcare}: Flagging uncertain diagnostic suggestions before they reach clinicians
        \item \textit{Legal}: Identifying potentially hallucinated legal citations or case references
        \item \textit{Education}: Warning students when AI-generated explanations may be unreliable
        \item \textit{Customer support}: Escalating uncertain responses to human agents
    \end{itemize}
    The finding that entropy-based uncertainty is domain-dependent (effective for knowledge questions, less so for reasoning) provides actionable guidance for deploying such systems: operators can calibrate confidence thresholds per domain.

    \item \textbf{Multilingual chatbot language control.} The language steering method from Study~2 addresses a concrete industry pain point: multilingual chatbots that unintentionally switch languages mid-response, degrading user experience. The method offers several advantages for industrial deployment:
    \begin{itemize}
        \item \textit{Zero training cost}: Steering requires only parallel translations for PCA calibration (as few as 50 sentence pairs), with no fine-tuning
        \item \textit{Negligible inference overhead}: One dot product and one vector subtraction per token
        \item \textit{No model modification}: Works with any decoder-only LLM via forward hooks
    \end{itemize}
    Potential applications include customer service platforms, content localization pipelines, and real-time translation assistants.

    \item \textbf{Cost-efficient domain adaptation.} The complexity-aware fine-tuning pipeline from Study~3 reduces the computational cost of adapting LLMs to specific domains by 79--82\% in terms of training data (tokens processed), while maintaining or improving accuracy. This has direct implications for:
    \begin{itemize}
        \item \textit{Enterprise AI}: Companies can fine-tune domain-specific models (medical, legal, financial) at a fraction of the current cost
        \item \textit{Resource-constrained settings}: Organizations with limited GPU budgets can achieve competitive performance
        \item \textit{Rapid iteration}: Faster training cycles enable more frequent model updates as domain knowledge evolves
    \end{itemize}
\end{enumerate}

\section*{Open-Source Contributions}

All three research projects have released their code and data publicly to facilitate reproducibility and further research:
\begin{enumerate}
    \item Uncertainty estimation pipeline: \texttt{github.com/LabARSS/question-complextiy-estimation}
    \item Complexity-aware fine-tuning: \texttt{github.com/LabARSS/complexity-aware-fine-tuning}
    \item Language steering: \texttt{github.com/fxlrnrpt/language-steering-in-latent-space}
\end{enumerate}

\section*{Start-up Potential}

The combination of the three methods developed in this thesis---uncertainty monitoring, language control, and efficient fine-tuning---forms a coherent product offering for an ``LLM reliability and control'' platform. Such a platform could provide:
\begin{itemize}
    \item Automated confidence scoring for LLM outputs (per-domain calibrated)
    \item Language consistency enforcement for multilingual deployments
    \item Efficient domain adaptation services with reduced compute requirements
\end{itemize}
All methods operate on standard transformer architectures without architectural modifications, making them compatible with the most widely deployed open-source LLMs.
